\section{CourseBoxes.sty}\label{sec:ref-courseboxes}

Theorem-like boxes in two families, both built on amsthm heads inside
tcolorbox frames:

\begin{itemize}
  \item \textbf{Text results} -- things that live in the course text and
    carry the number the text gives them. The number is \emph{declared}
    (typed at the call site), not counted, so the machinery is
    numbering-scheme-agnostic; \cs{label}/\cs{ref} record the typed number.
  \item \textbf{Asides} -- the instructor's pedagogical interjections.
    Unnumbered, and deliberately not \cs{label}-able: a stray label prints
    \texttt{[unlabelled aside]} -- loud -- rather than silently grabbing
    the section number.
\end{itemize}

Which family a name belongs to is a per-\emph{text} call made in the text
layer, not here: ``Key Idea'' is a numbered result in APEX and a plain
aside in a course whose text does not number it. \textbf{The factory you
declare it with is the decision} -- including whether it reaches a table of
contents (section~\ref{sec:ref-lecturenotes}), since results call the
contents hook and asides never do.

Load-order contract: amsthm must already be loaded, which Typefaces'
newtxmath owns -- so CourseBoxes comes \emph{after} Typefaces (enforced by
a \cs{PackageError}); hyperref (for the boxes' link anchors) arrives via
Hyperlinks in time for body use. Section~\ref{sec:load-order} has the whole
ordering story.

\subsection*{Shipped vocabulary}

Universal text results (every textbook has them):
\env{theorem}, \env{proposition}, \env{lemma}, \env{corollary} (italic
body); \env{definition}, \env{example}, \env{remark} (upright body). Call
site: \texttt{\char`\\begin\{theorem\}\{2.1.3\}[Power Rule]}.

\env{example} is the one text result that is \textbf{not boxed}, and
deliberately: the same word reaches the page two ways -- written here, or
as an imported problem shown as ``Example'' (section~\ref{sec:lectures}) --
and boxing one but not the other made two identical-sounding things look
unrelated. The imported route cannot be boxed safely (a problem carries
hints, solutions and \cs{workspace}), so neither is. The L-rule marks the
text's \emph{results}; a worked example is not one.

What then distinguishes an example \emph{from the text} from one of your
own is the \textbf{number}, not the frame --- which is what the starred
form below is for.

Shared asides (the author's voice, consistent course to course):
\env{caution} (red, warning glyph), \env{takeaway} (blue),
\env{yourturn} (blue, pen glyph), \env{prepwork} (blue, book glyph), and
\env{tip}\texttt{[surviving$\vert$thriving]} (orange; the optional tier
appends ``for the \emph{surviving/thriving}'' to the head).

Text-\emph{specific} vocabulary (Key Idea, Axiom, \dots) deliberately does
not ship here -- it is declared in the course's text layer,
\file{.textbook/textbook.tex}, with the factories below
(section~\ref{sec:extending}).

\subsection*{The factories}

\begin{center}
\begin{tabularx}{\linewidth}{@{}l Y@{}}
\cs{NewTextResult}\texttt{[head]\{env\}\{Name\}\{box\}} & declare a
  numbered text result; call site
  \texttt{\char`\\begin\{env\}\{num\}[name]}. Default head \env{cbdefn}
  (upright); pass \env{cbthm} for italic. An \emph{empty} box argument
  means no box at all (\env{example} uses this) -- not a blank frame:
  tcolorbox is never invoked, so such a result cannot affect pagination.\\
\cs{NewAside}\texttt{[head]\{env\}\{Name\}\{box\}} & declare an unnumbered
  aside. Default head \env{cbaside} (black sans).\\
\end{tabularx}
\end{center}

Every \cs{NewTextResult} also declares an \textbf{unnumbered twin},
\env{env*}, taking only the optional name:
\texttt{\char`\\begin\{example*\}[Cubic]} heads ``Example (Cubic)'' and
\texttt{\char`\\begin\{example*\}} just ``Example''. It shares the plain
form's head style and box, so the two cannot drift apart.

The number is mandatory in the plain form on purpose -- a declared number
is the norm -- so the star is the visible way to say ``this one has none''.
A starred result sets \texttt{[unnumbered \emph{Name}]} as its label text,
the same loud-sentinel trick asides use: without it an unnumbered result
would still be \cs{label}-able and \cs{ref} would expand to \emph{nothing},
leaving ``by Theorem~'' and a blank space in the PDF.

\subsection*{Restyling knobs}

Named once, referenced everywhere -- each is a one-line edit:

\begin{center}
\begin{tabularx}{\linewidth}{@{}l Y@{}}
tcolorbox styles & \env{cb/textresult} (quiet grey L-rule),
  \env{cb/note} (blue box), \env{cb/caution} (red), \env{cb/tip}
  (orange). An empty style is the ``no box'' case.\\
\cs{CourseBoxTocEntry}\texttt{\{Name\}\{num\}\{name\}} & called once per
  text result (never by an aside), a no-op here. A class that has a table
  of contents renews it; LectureNotes does
  (section~\ref{sec:ref-lecturenotes}).\\
amsthm head styles & \env{cbthm}, \env{cbdefn}, \env{cbaside},
  \env{cbnote}, \env{cbnotecolon} (colon punctuation), \env{cbcaution},
  \env{cbtip}. Head punctuation is the style's sixth argument.\\
colours & \env{ThmColor} (grey), \env{NoteColor} (royal blue),
  \env{CautionColor} (red), \env{TipColor} / \env{TipTextColor}
  (orange; the darker text value keeps WCAG contrast). Provisional
  pending a colour-blind-safe palette pass; the head glyphs
  (fontawesome, monochrome) are the colour-independent channel.\\
\end{tabularx}
\end{center}
